% ==========================
% General Conclusion
% ==========================
\CustomConclusion{\ThesisHeadingText{conclusion}}
\label{ch:conclusion}

This thesis presented a detailed theoretical and numerical study of exact solutions of the Einstein field equations, with particular emphasis on the Schwarzschild and Kerr metrics, in order to analyze their geodesic structure, horizon properties, and observable consequences for the behavior of matter, light, and gravitational radiation in strong gravitational fields.

Chapter~\ref{ch:chap1} recalled the geometric foundations of general relativity and derived the Einstein field equations of Eq.~\eqref{eq:efe} together with the Schwarzschild solution of Eq.~\eqref{eq:schwarzschild}, establishing the reference exact solution and the notion of the event horizon used throughout the remainder of the thesis \cite{einstein1916,schwarzschild1916}.

Chapter~\ref{ch:chap2} extended this framework to rotating compact objects through the Kerr metric of Eq.~\eqref{eq:kerr}, and showed, using the separability of the geodesic equations established by Carter \cite{carter1968}, how the innermost stable circular orbit depends on the spin parameter $a$, ranging from $3r_s$ in the Schwarzschild limit down to $0.5\,r_s$ for a maximally rotating black hole \cite{bardeen1972}.

Chapter~\ref{ch:chap3} implemented a numerical ray-tracing procedure to integrate null geodesics in both geometries, allowing the computation of gravitational lensing effects and the black hole shadow boundary set by the critical impact parameter of Eq.~\eqref{eq:critical-impact}. The resulting shadow sizes were found to be in good agreement with the Event Horizon Telescope observations of the supermassive black hole M87*, providing direct observational support for the strong-field predictions of the Kerr geometry \cite{eht2019}.

Chapter~\ref{ch:chap4} addressed the dynamical response of these spacetimes to small perturbations, introducing the quasinormal mode spectrum governed by the Regge--Wheeler and Teukolsky equations \cite{regge1957,teukolsky1973}. The dependence of this spectrum on the mass and spin of the remnant, combined with the no-hair theorem, was shown to underlie black hole spectroscopy, whose application to the ringdown of the gravitational-wave event GW150914 offers an independent dynamical confirmation of the Kerr hypothesis \cite{ligo2016,isi2019}.

Taken together, these results illustrate how a single family of exact solutions of the Einstein field equations accounts, in a unified manner, for phenomena as diverse as the orbital dynamics of accretion disks, the horizon-scale imaging of supermassive black holes, and the ringdown signals of compact binary mergers. The consistency between the analytical predictions derived in Chapters~\ref{ch:chap1}--\ref{ch:chap2} and the numerical and observational results of Chapters~\ref{ch:chap3}--\ref{ch:chap4} reinforces the Kerr metric's status as the standard model of astrophysical black holes.

Several directions naturally extend this work. On the theoretical side, the inclusion of an electric charge or a non-zero cosmological constant would lead to the more general Kerr--Newman and Kerr--de Sitter solutions, allowing the present analysis to be revisited in a broader class of exact solutions \cite{newman1965}. On the numerical side, extending the ray-tracing framework of Chapter~\ref{ch:chap3} to non-equatorial observers and to realistic accretion disk emission models would allow a more direct comparison with future high-resolution imaging campaigns. Finally, on the observational side, the increasing sensitivity of gravitational-wave detectors is expected to enable the measurement of multiple overtones in a single ringdown signal, strengthening the black hole spectroscopy tests introduced in Chapter~\ref{ch:chap4} and providing an increasingly stringent test of the no-hair theorem \cite{berti2009}.

% ==========================
% Last Line Shape 
% ==========================
\perfectdivider